
\section{Introduction}

Dynamical quantum systems involving both unitary evolution and stochastic wavefunction collapse from measurements provide a rich setting for new types of non-equilibrium matter~\cite{PotterVasseur2022,FisherKhemaniNahumVijay2023}.
In particular, the discovery of measurement-induced phase
transitions sparked rapid progress in our understanding of non-equilibrium physics~\cite{LiChenFisher2018,SkinnerRuhmanNahum2019,ChanNandkishorePretkoSmith2019,LiChenFisher2019,GullansHuse2020}. The relevant entanglenment dynamics phenonemology is further enriched by global symmetries~\cite{SangHsieh2021,BaoChoiAltman2021,LavasaniAlaviradBarkeshli2021,AgrawalEtAl2022,MajidyEtAl2023}, reshaped by unitary feedback conditioned on measurement outcomes~\cite{IadecolaEtAl2023,RavindranathEtAl2023,PiroliEtAl2023,SierantTurkeshi2023,ODeaEtAl2024}, and qualitatively modified by decoherence~\cite{WeinsteinBaoAltman2022,LiClaassen2023,LiSangHsieh2023,LiuLiZhangJian2024,IvakiOjanenMoghaddam2025}.

In this work, we seek to better understand the role of topology in these dynamical systems. In equilibrium, topological phases, such as Chern insulators, are a staple in
the study of quantum matter. For free fermion systems, topological phases are organized by symmetry and dimensionality into the tenfold way~\cite{AltlandZirnbauer1997,SchnyderRyuFurusakiLudwig2008,Kitaev2009,RyuSchnyderFurusakiLudwig2010,ChiuTeoSchnyderRyu2016}.
Topological phases cannot be smoothly deformed into a topologically distinct phase without closing the energy gap or breaking the relevant symmetry of the associated parent Hamiltonian. At the interface between distinct bulk topological phases, boundary-localized, gapless modes robust against symmetry-preserving
perturbations emerge~\cite{Halperin1982,Hatsugai1993,TeoKane2010}, and this bulk-boundary correspondence is a cornerstone of both experimental and theoretical condensed matter physics~\cite{HasanKane2010,QiZhang2011}. A key feature of the topological phases are their robustness to spatial disorder. In particular, the bulk invariant survives gap-and-symmetry-preserving spatial disorder, enabling the topologically-protected boundary modes which can
evade Anderson localization in the bulk, though this is specific to dimensionality~\cite{Halperin1982,EversMirlin2008}.

Monitored dynamical quantum systems provide a new arena in which to study this robustness, since the randomness of measurement outcomes acts as \textit{spatiotemporal} disorder. Unlike typical disorder, however, Born-rule conditioned measurement outcomes constitute temporal disorder that is correlated with the state itself, leading to distinct universality~\cite{JianShapourianBauerLudwig2023,FavaEtAl2023}. 

%In particular, this matters for gapless boundaries in monitored dynamical settings, since nonchiral free fermions in one dimension generically lose their critical correlations and remain area-law entangled at long time scales~\cite{PoboikoEtAl2023}, so it is not clear, a priori, that 1D critical states can be sustained under measurement-induced randomness.

Recently, the role of topology in monitored quantum dynamics has advanced notably, particularly for free fermions. Concrete $1+1$D models with topologically distinct area-law phases have been identified~\cite{KellsMeidanRomito2023,BehrendsVennBeri2024,PanShapourianJian2025,Oshima2025, yang2026decoding}, while Refs.~\cite{BhuiyanPanJian2026,XiaoKawabata2026} develop topological classifications of free-fermion monitored dynamics based on the AZ symmetry classification.

There has also been notable progress in the physics of boundary modes in topological dynamical phases.  Refs.~\cite{PanShapourianJian2025, Oshima2025} both identified dynamical topological modes localized on $0+1$D domain walls between $1+1$D topological dynamical phases. Ref.~\cite{XiaoKawabata2026}
established an analog of the bulk-boundary correspondence in which non-trivial bulk topology corresponds to the gapless modes in the Lyapunov spectrum of the dynamics. Their $2+1$D construction, however, postselects the measurement outcomes, so it is unclear if chiral boundary modes survive in a stochastic, monitored setting. Whether chiral, gapless boundary modes can be realized under genuine Born-rule dynamics is of interest to this work.

%\claudeadd{The question is also natural from the perspective of state preparation. A single chiral branch is anomalous: it can only exist at the boundary of a two-dimensional region with nonzero Chern number. Preparing it therefore requires preparing a chiral bulk, which strictly local finite-depth unitary circuits cannot do~\cite{LuLessaKimHsieh2022}, and which finite-range dissipative dynamics cannot stabilize as a pure steady state~\cite{BardynEtAl2013,BudichZollerDiehl2015,Goldstein2019}. Adaptive circuits that combine measurements with classical feedforward circumvent some of these constraints. Reference~\cite{LuLessaKimHsieh2022} prepares chiral topological order in constant depth, given a chiral free-fermion resource state, and one-dimensional critical states in logarithmic depth using nonlocal classical communication. Here we ask whether a much simpler ingredient, repeated finite-range measurements with purely local feedback, can prepare chiral critical boundaries directly. We find that it can, but at the level of individual trajectories: each measurement record yields a pure state whose walls are critical and chiral, whereas the outcome-averaged state is short-range correlated.}


%In this paper, we consider a 2+1D adaptive quantum circuit consisting of finite-range occupation measurements, local unitary feedback, and ancilla modes—a variant of the model introduced in Ref.~\cite{BhuiyanPanJian2026}. The circuit steers toward the topological phase by measuring truncated, overcomplete Wannier (OW) modes and swapping in fresh ancilla modes with the target occupancy whenever an undesired physical-system occupancy is observed. Truncation mixes the two bands, rendering the finite-range OW stabilizers frustrated, with no common dark state. This frustration is unavoidable: compactly supported single-particle functions cannot both lie within a Chern band and span it at every momentum~\cite{DubailRead2015}. Instead, the dynamics approach a disordered ensemble of Chern insulator states.

 
%By tracing out ancilla degrees of freedom, we show that unitary feedback becomes equivalent to fermion-counting operations on the physical system, consisting of particle gain and loss operations in the OW basis that correct undesired measurement outcomes. This simplified model will be the main system of interest in this work.

%We introduce programmable domain walls by making the circuit parameter $\alpha$, which controls the topological phase, spatially dependent and truncating OW mode supports at the interfaces. We first verify that the adaptive circuit prepares a topological bulk within a finite slab, then show algebraically slow purification when the slab is tuned to the topological phase. We connect this slow purification to a finite-time Lyapunov gap that decreases with system size and to slow modes localized near the boundaries and extended along them. We then study pure states along individual quantum trajectories, finding algebraically decaying correlations and logarithmic scaling of both entanglement and charge variance, consistent with equilibrium expectations. The trajectory-resolved spectral densities provide complementary evidence. We further examine the convergence of the entropy coefficient, resolve the contribution from each wall, and use modular evolution to identify opposite chiralities.  Finally, we examine the trajectory-averaged quantum channel, finding short-range two-point correlations and a finite relaxation gap over the sizes studied. These results indicate that the boundary criticality resides in the conditioned trajectory ensemble and is absent in the outcome-averaged state.

%Reference~\cite{LuLessaKimHsieh2022} prepares chiral topological order in constant depth, given a chiral free-fermion resource state, and one-dimensional critical states in logarithmic depth using nonlocal classical communication. Here we ask whether a much simpler ingredient, repeated finite-range measurements with purely local feedback, can prepare chiral critical boundaries directly. We find that it can, but at the level of individual trajectories: each measurement record yields a pure state whose walls are critical and chiral, whereas the outcome-averaged state is short-range correlated.

%Together, these results provide evidence for critical chiral boundary state ensembles in finite-range, postselection-free adaptive dynamics. We do not claim that the trajectory-conditioned boundary modes coincide with the ground state of an equilibrium chiral Luttinger liquid; rather, our work concerns universal features of trajectory-conditioned correlations, entanglement, and chirality, diagnosed at the ensemble level, and provides a post-selection-free protocol to prepare topologically-protected chiral critical boundary modes.

The question is also natural from the perspective of state preparation. A single chiral branch is anomalous, as it can only exist at the boundary of a two-dimensional region with nonzero Chern number. Preparing it therefore requires preparing a chiral bulk, which strictly local finite-depth unitary circuits cannot do~\cite{LuLessaKimHsieh2022}, and which finite-range dissipative dynamics cannot stabilize as a pure steady state~\cite{BardynEtAl2013,BudichZollerDiehl2015,Goldstein2019}. Adaptive circuits that combine measurements with classical feedforward circumvent some of these constraints. Reference~\cite{LuLessaKimHsieh2022} prepares chiral topological order in constant depth, given a chiral free-fermion resource state, and one-dimensional critical states in logarithmic depth using nonlocal classical communication. Here we ask whether a much simpler ingredient, repeated finite-range measurements with purely local feedback, can prepare chiral critical boundaries directly. We find that it can, but at the level of individual trajectories, where each measurement record yields a pure state with boundary-localized critical, chiral modes, whereas the trajectory-averaged state shows short-range correlated boundary-localized mixed states.

In this paper, we consider a 2+1D adaptive quantum circuit consisting of finite-range occupation measurements, local unitary feedback, and ancilla modes---a variant of the model introduced in Ref.~\cite{BhuiyanPanJian2026}. The circuit steers toward the topological phase by measuring truncated, overcomplete Wannier (OW) modes and swapping in fresh ancilla modes with the target occupancy whenever an undesired physical-system occupancy is observed. Truncation mixes the two bands, rendering the finite-range OW stabilizers frustrated, with no common dark state. This frustration is unavoidable, since compactly supported single-particle functions cannot both lie within a Chern band and span it at every momentum~\cite{DubailRead2015}. Instead, the dynamics approach a disordered ensemble of Chern insulator states.

We introduce programmable domain walls by making the circuit parameter $\alpha$, which controls the topological phase, spatially dependent and truncating OW mode supports at the interfaces. We first verify that the adaptive circuit prepares a topological bulk within a finite slab, then show algebraically slow purification when the slab is tuned to the topological phase. We connect this slow purification to a finite-time Lyapunov gap that decreases with system size and to slow modes localized near the boundaries and extended along them. We then study pure states along individual quantum trajectories, finding algebraically decaying correlations and logarithmic scaling of both entanglement and charge variance, with coefficients close to those of free chiral fermions. The trajectory-resolved spectral densities provide complementary evidence. We further examine the convergence of the entropy coefficient, resolve the contribution from each wall, and use modular evolution to identify opposite chiralities. Finally, we examine the trajectory-averaged quantum channel, finding short-range two-point correlations and a finite relaxation gap over the sizes studied.

Together, these results provide evidence for critical chiral boundary state ensembles in finite-range, postselection-free adaptive dynamics. We do not claim that the trajectory-conditioned boundary modes coincide with the ground state of an equilibrium chiral Luttinger liquid. Rather, our work concerns universal features of correlations, entanglement, and chirality, characterized at the ensemble level of Born-rule conditioned trajectories.